%document
\documentclass[10pt]{beamer}
%theme
\usetheme{CambridgeUS}
% packages
\usepackage{color}
\usepackage{listings}
\usepackage[ngerman]{babel}
\usepackage[utf8]{inputenc}
\usepackage{multicol}


% color definitions
\definecolor{mygreen}{rgb}{0,0.6,0}
\definecolor{mygray}{rgb}{0.5,0.5,0.5}
\definecolor{mymauve}{rgb}{0.58,0,0.82}

\lstset{
    backgroundcolor=\color{white},
    % choose the background color;
    % you must add \usepackage{color} or \usepackage{xcolor}
    basicstyle=\footnotesize\ttfamily,
    % the size of the fonts that are used for the code
    breakatwhitespace=false,
    % sets if automatic breaks should only happen at whitespace
    breaklines=true,                 % sets automatic line breaking
    captionpos=b,                    % sets the caption-position to bottom
    commentstyle=\color{mygreen},    % comment style
    % deletekeywords={...},
    % if you want to delete keywords from the given language
    extendedchars=true,
    % lets you use non-ASCII characters;
    % for 8-bits encodings only, does not work with UTF-8
    frame=single,                    % adds a frame around the code
    keepspaces=true,
    % keeps spaces in text,
    % useful for keeping indentation of code
    % (possibly needs columns=flexible)
    keywordstyle=\color{blue},       % keyword style
    % morekeywords={*,...},
    % if you want to add more keywords to the set
    numbers=left,
    % where to put the line-numbers; possible values are (none, left, right)
    numbersep=5pt,
    % how far the line-numbers are from the code
    numberstyle=\tiny\color{mygray},
    % the style that is used for the line-numbers
    rulecolor=\color{black},
    % if not set, the frame-color may be changed on line-breaks
    % within not-black text (e.g. comments (green here))
    stepnumber=1,
    % the step between two line-numbers.
    % If it's 1, each line will be numbered
    stringstyle=\color{mymauve},     % string literal style
    tabsize=4,                       % sets default tabsize to 4 spaces
    % show the filename of files included with \lstinputlisting;
    % also try caption instead of title
    language = Java,
	showspaces = false,
	showtabs = false,
	showstringspaces = false,
	escapechar = ,
}

\def\ContinueLineNumber{\lstset{firstnumber=last}}
\def\StartLineAt#1{\lstset{firstnumber=#1}}
\let\numberLineAt\StartLineAt



\newcommand{\codeline}[1]{
	\alert{\texttt{#1}}
}

% This Document contains the information about this course.

% Authors of the slides
\author{Vincent Gerber, Tilman Hinnerichs}

% Name of the Course
\institute{Java Kurs}

% Fancy Logo
\titlegraphic{\hfill\includegraphics[height=1.25cm]{../templates/fsr_logo_cropped}}


\usepackage{qtree}

\title{Java}
\subtitle{Java files IO}
\date{\today}

\usepackage{color}

\begin{document}
\begin{frame}
	\titlepage
\end{frame}

\begin{frame}{Overview}
	\tableofcontents
\end{frame}

\section{How to use streams}
\begin{frame}{Streams}
	A stream can be defined as a sequence of data. We got two kinds of streams in Java:
	\vfill
	\begin{tabular}{rl}
		InPutStream&is used to read data from a source\\
		OutPutStream&is used for writing data to a destination
	\end{tabular}
	\vfill
	Many implementations can be found in the module \textbf{java.io}.
\end{frame}

\begin{frame}[fragile]{File Streams}
	Byte streams used to perform input and ouput of 8-bit bytes. The most common streams for file reading and writing are \textbf{FileInputStream} and \textbf{FileOutputStream}:\\
	
\end{frame}

\begin{frame}[fragile]{How to use file streams}
	We would like to copy a file "trateotu.txt":
	\begin{lstlisting}
	import java.io.*
	... main(String args[]) throws IOException {
		FileInputStream in = null;
		FileOutputStream out = null;
		
		try {
			in = new FileInputStream("trateotu-input.txt");
			out = new FileOutputStream("trateotu-output.txt");
			
			int c;
			while ((c = in.read()) != -1) {
				out.write(c);
			}
		}
		finally {
			if (in != null) {
				in.close();
			}
			if (out != null) {
				out.close();
			}
		}
	}
	\end{lstlisting}
\end{frame}


\begin{frame}{Other useful methods for InputStream}
	\begingroup
	\def\arraystretch{1.5}
	\begin{tabular}{rp{8cm}}
		public void close()&releases any system resources associated with the file\\
		protected void finalize()&cleans up connection to the file, closes stream when there are no more references\\
		public int read(int r)&this method reads the specified byte of data\\
		public int available()& returns the amount of bytes which can be read from the input stream
	\end{tabular}
	\endgroup
\end{frame}

\begin{frame}{Othe useful methods for OutputStream}
	
	\begingroup
	\def\arraystretch{1.5}
	\begin{tabular}{rp{8cm}}
		public void close()&closes output stream, releases any resources\\
		public void finalize()&see InputStream\\
		public void write(int w)&writes the specified byte
	\end{tabular}
	\endgroup
\end{frame}

\begin{frame}[fragile]{Character Streams}
	Instead of the previously used Byte streams, we can also use \textbf{Character streams}. \textbf{Character streams} are able to perform input and output for 16-bit unicode:
	
	\begin{lstlisting}
		in = new FileReader("ltuae-input.txt);
		out = new FileReader("ltuae-output.txt);
	\end{lstlisting}
\end{frame}

\begin{frame}{User interaction}
	Now knowing about streams we can also perform command line input/output. Java provides the following three streams, given by your OS:\\
	\vfill
	\begingroup
	\def\arraystretch{1.5}
	\begin{tabular}{r p{8cm}}
		Standard Input&is used to feed data into you program, is your standard input stream, represented as \textbf{System.in}\\
		Standard Output & is used to output the data produced by the user's program, represented as \textbf{System.out}\\
		Standard Error&is used to output the error data produced by our program, \textbf{System.err}
	\end{tabular}
	\endgroup
\end{frame}

\begin{frame}[fragile]{First Console IO program}
	We will now read standard input stream until the user types a "q":\\
	\begin{lstlisting}
		InputStreamReader cin = null;
		
		try {
			cin = new InputStreamReader(System.in);
			System.out.println("Enter characters, 'q' to quit.");
			char c;
			do {
				c = (char) cin.read();
				System.out.print(c);
			} while(c != 'q');
		}
		finally {
			if (cin != null) {
				cin.close();
			}
		}
	\end{lstlisting}
\end{frame}

\end{document}